\newpage
\section{\textbf{Dimensionamento e especificação dos TCs}}

% ===================== DIMENSIONAMENTO E ESPECIFICAÇÃO DOS TCs =====================

Os transformadores de corrente (TCs) reproduzem, em seu secundário, uma imagem reduzida das correntes primárias, para que os relés de proteção possam medi-las. Para que a proteção opere corretamente, o TC deve suportar a corrente de carga em regime permanente e reproduzir as correntes de curto-circuito sem saturar, de modo que o relé receba um sinal fiel à corrente de falta \cite{mamede2013}.

Nesta seção são dimensionados os TCs de proteção de sobrecorrente da SE Campus do Pici, identificados no diagrama unifilar da Figura~\ref{fig:diagrama_unifilar_protecao}: os TCs de fase do vão de entrada de linha (12L1), dos disjuntores de alta tensão dos transformadores (12T1 e 12T2), dos disjuntores gerais de média tensão (11T1 e 11T2) e das saídas dos alimentadores (21L1, 21L2 e 21L3), além do TC de neutro dos transformadores, utilizado pela função 51G. Os TCs associados à proteção diferencial dos transformadores são tratados na seção referente à função 87T.

No dimensionamento, considerou-se a operação da subestação com apenas um transformador energizado, permanecendo o outro como reserva. Os TCs dos dois transformadores recebem a mesma especificação, de modo que qualquer um deles possa entrar em operação sem alteração nos ajustes da proteção.

\subsection{Dados de projeto}

Os TCs disponíveis para os setores de alta e média tensão da subestação são apresentados na Tabela~\ref{tab:tc_disponiveis}. Todos são do tipo multirrelação, com fator de sobrecorrente $FS = 20$, fator térmico $FT = 1{,}2$ e classe de exatidão 10B200, isto é, erro máximo de 10\% para uma tensão secundária de até 200~V, desenvolvida com 20 vezes a corrente nominal secundária (100~A) circulando na carga nominal de 2~$\Omega$.

\begin{table}[htbp]
    \centering
    \caption{Transformadores de corrente disponíveis para a SE Campus do Pici.}
    \label{tab:tc_disponiveis}
    \begin{tabular}{ccccc}
        \hline
        \textbf{Nível de tensão} & \textbf{Relação nominal (A)} & \textbf{FS} & \textbf{FT} & \textbf{Exatidão} \\
        \hline
        69 kV   & 200/400/600-5          & 20 & 1,2 & 10B200 \\
        \hline
        \multirow{3}{*}{13,8 kV} & 800/1000/1200/1400-5 & 20 & 1,2 & 10B200 \\
                                 & 200/400-5            & 20 & 1,2 & 10B200 \\
                                 & 100/200/400/600/800-5 & 20 & 1,2 & 10B200 \\
        \hline
    \end{tabular}
    \par\vspace{2pt}
    {\small Fonte: \cite{sampaio2026roteiro}.}
\end{table}

A carga conectada ao secundário de cada TC é formada pelos condutores que o interligam ao relé na casa de comando e pela carga do próprio relé. Os dados utilizados estão na Tabela~\ref{tab:dados_carga_tc}.

\begin{table}[htbp]
    \centering
    \caption{Dados dos condutores secundários e do relé.}
    \label{tab:dados_carga_tc}
    \begin{tabular}{lcc}
        \hline
        \textbf{Dado} & \textbf{Pátio de 69 kV} & \textbf{Pátio de 13,8 kV} \\
        \hline
        Comprimento dos condutores ($L$)  & 100 m & 80 m \\
        Impedância dos condutores ($Z_c$) & $(5{,}52 + j0{,}13)~\Omega$/km & $(5{,}52 + j0{,}13)~\Omega$/km \\
        Carga do relé ($S_r$)             & 0,2 VA & 0,2 VA \\
        \hline
    \end{tabular}
    \par\vspace{2pt}
    {\small Fonte: \cite{sampaio2026roteiro}.}
\end{table}

A impedância do relé, referida à corrente nominal secundária de 5~A, é dada por:

\begin{equation}
    Z_r = \frac{S_r}{I_{sec}^2} = \frac{0{,}2}{5^2} = 0{,}008~\Omega
\end{equation}

A impedância total no secundário do TC é calculada pela Equação~\eqref{eq:z_total}. Considera-se o percurso de ida e volta dos condutores ($2L$), que corresponde à condição mais severa, verificada nas faltas fase-terra, quando a corrente retorna pelo condutor de neutro do circuito secundário. A resistência interna do enrolamento secundário do TC foi desprezada, por não ter sido fornecida.

\begin{equation}
    Z_{TC} = 2 \cdot L \cdot Z_c + Z_r
    \label{eq:z_total}
\end{equation}

\begin{itemize}
    \item[--] Pátio de 69 kV: $Z_{TC} = 2 \cdot 0{,}100 \cdot (5{,}52 + j0{,}13) + 0{,}008 = 1{,}112 + j0{,}026~\Omega \Rightarrow |Z_{TC}| = 1{,}1123~\Omega$;
    \item[--] Pátio de 13,8 kV: $Z_{TC} = 2 \cdot 0{,}080 \cdot (5{,}52 + j0{,}13) + 0{,}008 = 0{,}8912 + j0{,}0208~\Omega \Rightarrow |Z_{TC}| = 0{,}8914~\Omega$.
\end{itemize}

As correntes máximas de carga e de curto-circuito consideradas para cada TC estão resumidas na Tabela~\ref{tab:correntes_tc}, adotando-se os seguintes critérios:

\begin{itemize}
    \item[--] \textbf{12L1}: com apenas um transformador energizado, a corrente de carga corresponde à corrente nominal ONAF desse transformador em 69~kV (52,3~A), e a corrente de curto-circuito máxima é a falta trifásica no barramento de 69~kV (15.683,2~A);
    \item[--] \textbf{12T1 e 12T2}: a corrente de carga é a corrente nominal ONAF de cada transformador em 69~kV (52,3~A). A corrente de curto-circuito máxima ocorre para uma falta nos terminais de alta tensão do transformador, alimentada pelo barramento de 69~kV, e é igual à corrente de falta trifásica nessa barra (15.683,2~A);
    \item[--] \textbf{11T1 e 11T2}: a corrente de carga é a corrente nominal ONAF do transformador em 13,8~kV (261,5~A). Como apenas um transformador está em operação, toda a corrente de falta no barramento de 13,8~kV circula pelo seu disjuntor geral, e a corrente de curto-circuito máxima é a falta fase-terra nessa barra (2.918,4~A);
    \item[--] \textbf{21L1, 21L2 e 21L3}: como os relés dos alimentadores devem ter o mesmo ajuste \cite{sampaio2026roteiro}, adotou-se a maior corrente de carga registrada entre os alimentadores (100~A). A corrente de curto-circuito máxima é a falta fase-terra no barramento de 13,8~kV (2.918,4~A);
    \item[--] \textbf{Neutro dos transformadores}: em regime normal, o TC de neutro conduz apenas a corrente de desequilíbrio. Sua corrente máxima é a corrente residual ($3I_0$) da falta fase-terra no barramento de 13,8~kV, que retorna integralmente pelo neutro do transformador em operação (2.918,4~A).
\end{itemize}

\begin{table}[htbp]
    \centering
    \caption{Correntes máximas de carga e de curto-circuito consideradas no dimensionamento dos TCs.}
    \label{tab:correntes_tc}
    \begin{tabular}{lccc}
        \hline
        \textbf{TC} & \textbf{Tensão (kV)} & \textbf{$I_{carga,máx}$ (A)} & \textbf{$I_{cc,máx}$ (A)} \\
        \hline
        12L1                 & 69   & 52,3  & 15.683,2 \\
        12T1 e 12T2          & 69   & 52,3  & 15.683,2 \\
        11T1 e 11T2          & 13,8 & 261,5 & 2.918,4 \\
        21L1, 21L2 e 21L3    & 13,8 & 100,0 & 2.918,4 \\
        Neutro 02T1 e 02T2   & 13,8 & --    & 2.918,4 \\
        \hline
    \end{tabular}
    \par\vspace{2pt}
    {\small Fonte: elaborado pelos autores.}
\end{table}

\subsection{Critérios de dimensionamento}

Os TCs foram dimensionados segundo os quatro critérios estabelecidos em \cite{sampaio2026roteiro}, descritos a seguir.

\textbf{a) Critério da corrente nominal}: a corrente primária nominal do TC ($I_p$) deve ser maior ou igual à corrente máxima de carga. Verifica-se também se a corrente admissível em regime permanente, dada pelo fator térmico, atende à carga máxima:

\begin{equation}
    I_p \geq I_{carga,máx} \qquad \text{e} \qquad FT \cdot I_p \geq I_{carga,máx}
    \label{eq:crit_nominal}
\end{equation}

\textbf{b) Critério do curto-circuito}: a corrente primária nominal do TC multiplicada pelo fator de sobrecorrente deve ser maior ou igual à corrente máxima de curto-circuito:

\begin{equation}
    FS \cdot I_p \geq I_{cc,máx} \quad \Rightarrow \quad I_p \geq \frac{I_{cc,máx}}{20}
    \label{eq:crit_cc}
\end{equation}

\textbf{c) Critério da tensão máxima no secundário}: a tensão desenvolvida no secundário do TC durante o curto-circuito máximo deve ser menor que a tensão nominal da classe de exatidão ($V_N = 200$~V):

\begin{equation}
    V_s = \frac{I_{cc,máx}}{RTC} \cdot |Z_{TC}| \leq V_N
    \label{eq:crit_tensao}
\end{equation}

\textbf{d) Critério de Stanley E. Zocholl}: verifica se o TC reproduz, sem saturar, a corrente de curto-circuito com a máxima componente contínua (assimetria total). O critério é expresso por:

\begin{equation}
    20 \geq \left(\frac{X}{R} + 1\right) \cdot I_f \cdot Z_b
    \label{eq:zocholl}
\end{equation}

\noindent em que $I_f$ é a corrente máxima de falta no secundário do TC em pu da corrente nominal secundária e $Z_b$ é a carga do TC em pu da carga nominal da classe de exatidão. Detalhando os termos:

\begin{equation}
    20 \geq \left(\frac{X}{R} + 1\right) \cdot \left(\frac{I_{cc,máx,sec}}{I_{TC,sec}}\right) \cdot \left(\frac{Z_{TC}}{\dfrac{V_{TC,máx}}{20 \cdot I_{TC,sec}}}\right)
    \label{eq:zocholl_detalhada}
\end{equation}

Para a classe 10B200, a carga nominal é $V_{TC,máx}/(20 \cdot I_{TC,sec}) = 200/(20 \cdot 5) = 2~\Omega$. Foi adotada a relação $X/R = 12{,}77$ \cite{sampaio2026roteiro}.

Entre as relações que atendem aos critérios a, b e c, foi escolhida a menor. Uma RTC menor resulta em correntes secundárias maiores para um mesmo valor primário, o que permite ajustar as funções de neutro (51N, 51NS e 51G) com valores de \textit{pick-up} reduzidos, dentro das faixas de ajuste do relé \cite{sel2013sel751}.

\subsection{Dimensionamento do TC do disjuntor de entrada 12L1}

\textbf{a) Critério da corrente nominal}

\begin{equation*}
    I_p \geq 52{,}3~\text{A}
\end{equation*}

Pelo critério da corrente nominal, a menor relação disponível (200-5) já seria suficiente.

\textbf{b) Critério do curto-circuito}

\begin{equation*}
    I_p \geq \frac{15.683{,}2}{20} = 784{,}2~\text{A}
\end{equation*}

A maior relação disponível para 69~kV é 600-5. Como $20 \cdot 600 = 12.000~\text{A} < 15.683{,}2~\text{A}$, o critério do curto-circuito não é atendido por nenhum dos TCs disponíveis para esse nível de tensão. Assim, adotou-se a maior relação, 600-5 ($RTC = 120$), e sua adequação foi verificada pelo critério da tensão máxima no secundário.

\textbf{c) Critério da tensão máxima no secundário}

\begin{equation*}
    V_s = \frac{15.683{,}2}{120} \cdot 1{,}1123 = 130{,}69 \cdot 1{,}1123 = 145{,}4~\text{V} \leq 200~\text{V}
\end{equation*}

O critério é atendido. Embora a corrente de falta corresponda a 26,1 vezes a corrente nominal do TC, a carga real conectada ao secundário (1,11~$\Omega$) é praticamente a metade da carga nominal da classe (2~$\Omega$). Por isso, o TC pode conduzir até $200/(5 \cdot 1{,}1123) = 35{,}96$ vezes sua corrente nominal sem ultrapassar a tensão de 200~V, ou seja, ele reproduz a corrente de falta simétrica sem saturar.

\textbf{d) Critério de Stanley E. Zocholl}

\begin{equation*}
    I_f = \frac{130{,}69}{5} = 26{,}14~\text{pu} \qquad Z_b = \frac{1{,}1123}{2} = 0{,}5562~\text{pu}
\end{equation*}

\begin{equation*}
    \left(12{,}77 + 1\right) \cdot 26{,}14 \cdot 0{,}5562 = 200{,}2 > 20
\end{equation*}

O critério de Zocholl não é atendido, o que indica que o TC pode saturar durante os primeiros ciclos de um curto-circuito com assimetria máxima.

\subsection{Dimensionamento dos demais TCs}

Os demais TCs foram dimensionados com o mesmo procedimento, e os resultados estão resumidos na Tabela~\ref{tab:resultados_tc}.

\begin{table}[htbp]
    \centering
    \caption{Resultados dos critérios de dimensionamento dos TCs.}
    \label{tab:resultados_tc}
    \resizebox{\textwidth}{!}{%
    \begin{tabular}{lccccccccc}
        \hline
        \multirow{2}{*}{\textbf{TC}} & \multirow{2}{*}{\textbf{RTC}} & \multicolumn{2}{c}{\textbf{a) Corrente nominal}} & \multicolumn{2}{c}{\textbf{b) Curto-circuito}} & \multicolumn{2}{c}{\textbf{c) Tensão no secundário}} & \multicolumn{2}{c}{\textbf{d) Zocholl}} \\
         & & $I_p \geq I_{carga}$ & Atende? & $I_{cc}/20$ (A) & Atende? & $V_s$ (V) & Atende? & Resultado & Atende? \\
        \hline
        12L1               & 600-5 (120) & $600 \geq 52{,}3$  & Sim & 784,2 & Não & 145,4 & Sim & 200,2 & Não \\
        12T1 e 12T2        & 600-5 (120) & $600 \geq 52{,}3$  & Sim & 784,2 & Não & 145,4 & Sim & 200,2 & Não \\
        11T1 e 11T2        & 400-5 (80)  & $400 \geq 261{,}5$ & Sim & 145,9 & Sim & 32,5  & Sim & 44,8  & Não \\
        21L1, 21L2 e 21L3  & 200-5 (40)  & $200 \geq 100{,}0$ & Sim & 145,9 & Sim & 65,0  & Sim & 89,6  & Não \\
        Neutro 02T1 e 02T2 & 200-5 (40)  & --                 & --  & 145,9 & Sim & 65,0  & Sim & 89,6  & Não \\
        \hline
    \end{tabular}}
    \par\vspace{2pt}
    {\small Fonte: elaborado pelos autores.}
\end{table}

Para os TCs 12T1 e 12T2, assim como para o TC 12L1, a corrente de carga de 52,3~A seria atendida pela relação 200-5. No entanto, como esses TCs estão sujeitos à mesma corrente de falta do TC 12L1, foi adotada a relação 600-5, que é a única que atende ao critério da tensão máxima no secundário. Com essa relação, uma corrente de \textit{pick-up} da função 51 de $1{,}25 \times 52{,}3 = 65{,}4$~A corresponde a 0,545~A no secundário, ainda acima do valor mínimo de ajuste do relé SEL-751 para TCs de 5~A (0,50~A) \cite{sel2013sel751}.

Para os TCs 11T1 e 11T2, a relação 400-5 é a menor que atende à corrente nominal ONAF do transformador. Essa relação também atende à soma das cargas máximas dos três alimentadores ($46 + 49{,}6 + 100 = 195{,}6$~A), suprida integralmente pelo transformador em operação. Seu limite térmico em regime permanente ($1{,}2 \times 400 = 480$~A) comporta ainda a sobrecarga admissível do transformador considerada no ajuste da função 51 (até $1{,}5 \times 261{,}5 = 392{,}2$~A).

Para os TCs dos alimentadores, a relação 200-5 atende à maior corrente de carga registrada (100~A) e ao critério do curto-circuito ($20 \times 200 = 4.000~\text{A} \geq 2.918{,}4~\text{A}$). Seu limite térmico em regime permanente ($1{,}2 \times 200 = 240$~A) é superior à soma das cargas máximas dos três alimentadores ($46 + 49{,}6 + 100 = 195{,}6$~A). Dessa forma, o TC suporta inclusive a transferência de toda a carga do campus para um único alimentador por meio dos pontos de interligação RP1 e RP2. Como os TCs são multirrelação, a derivação pode ser alterada para 400-5 caso a carga dos alimentadores cresça até valores próximos da capacidade dos condutores (266~A).

\subsection{Análise do critério de Zocholl}

Nenhum dos TCs atende ao critério de Zocholl. Isso indica que, para curtos-circuitos próximos à subestação com assimetria máxima, os TCs podem saturar nos primeiros ciclos da falta, enquanto a componente contínua da corrente não se extingue. Para a relação $X/R = 12{,}77$, a constante de tempo dessa componente é:

\begin{equation}
    \tau = \frac{X/R}{2 \pi f} = \frac{12{,}77}{2 \pi \cdot 60} = 33{,}9~\text{ms}
\end{equation}

\noindent ou seja, a componente contínua é praticamente extinta após cerca de $5\tau \approx 170$~ms. Assim, a saturação transitória tende a afetar pouco as funções temporizadas (51 e 51N), cujos tempos de operação para as correntes de falta em questão são da ordem de centenas de milissegundos, e afeta principalmente as funções instantâneas (50 e 50N), que podem ter a operação retardada nos primeiros ciclos da falta.

No setor de 13,8~kV, o critério seria atendido com a redução da carga secundária para valores de até 0,40~$\Omega$ (TCs 11T1 e 11T2) e 0,20~$\Omega$ (TCs dos alimentadores e de neutro), o que pode ser obtido com o aumento da seção dos condutores secundários ou a redução da distância entre os TCs e os relés. No setor de 69~kV, mesmo com a relação de 600-5, a carga secundária teria de ser inferior a 0,11~$\Omega$, valor menor que a impedância dos próprios condutores, o que exigiria TCs de classe de exatidão superior. Essas alternativas são discutidas na seção de proposição de melhorias.

\subsection{Especificação dos TCs}

A especificação final dos TCs de proteção de sobrecorrente da SE Campus do Pici é apresentada na Tabela~\ref{tab:especificacao_tc}.

\begin{table}[htbp]
    \centering
    \caption{Especificação dos TCs de proteção de sobrecorrente da SE Campus do Pici.}
    \label{tab:especificacao_tc}
    \resizebox{\textwidth}{!}{%
    \begin{tabular}{lccccccc}
        \hline
        \textbf{TC} & \textbf{Qtd.} & \textbf{Tensão (kV)} & \textbf{Relação nominal (A)} & \textbf{Derivação (RTC)} & \textbf{FS} & \textbf{FT} & \textbf{Exatidão} \\
        \hline
        12L1               & 3 & 69   & 200/400/600-5         & 600-5 (120) & 20 & 1,2 & 10B200 \\
        12T1 e 12T2        & 6 & 69   & 200/400/600-5         & 600-5 (120) & 20 & 1,2 & 10B200 \\
        11T1 e 11T2        & 6 & 13,8 & 200/400-5             & 400-5 (80)  & 20 & 1,2 & 10B200 \\
        21L1, 21L2 e 21L3  & 9 & 13,8 & 100/200/400/600/800-5 & 200-5 (40)  & 20 & 1,2 & 10B200 \\
        Neutro 02T1 e 02T2 & 2 & 13,8 & 100/200/400/600/800-5 & 200-5 (40)  & 20 & 1,2 & 10B200 \\
        \hline
    \end{tabular}}
    \par\vspace{2pt}
    {\small Fonte: elaborado pelos autores.}
\end{table}
